\begin{enumerate}
\item \textbf{Mensaje de arranque.} En la pantalla del propio MSX, \texttt{OCEAN SOFTWARE 1989 / DIHALT STUDIO 2026 / V9990 found !} durante unos dos segundos.
\item \textbf{Logo de DiHalt.} Se ve cuatro segundos y no se puede saltar.
\item \textbf{Portada.} La ilustración de la recreativa. Pulsa cualquier tecla o un botón del joystick para seguir.
\item \textbf{Pantalla de título.} \texttt{CHASE H.Q.} con \texttt{PRESS GEAR TO PLAY} parpadeando; debajo se alternan los créditos de la versión MSX y el copyright. Pulsa \textbf{MARCHA} (\key{ESPACIO} o el botón 1 del joystick) para empezar.
\item \textbf{Aviso.} La pantalla \texttt{CHASE H.Q. MONITORING SYSTEM}: Nancy, desde Chase H.Q., te cuenta por radio quién es el criminal, y el monitor muestra su coche. Pulsa \textbf{MARCHA} para saltarlo.
\item \textbf{La persecución.} Tras cada detención empieza el aviso de la fase siguiente. Hay cinco fases; después de la quinta viene el final y luego la pantalla de título.
\end{enumerate}

Entre una pantalla y otra la imagen se queda un momento en negro mientras se prepara la siguiente. Cuando se acaba la partida sale \texttt{GAME OVER} y vuelve la pantalla de título.
